\documentclass[a4paper]{article}
\usepackage{graphicx}
\usepackage{amsmath,amssymb}
\usepackage{hyperref}
\usepackage{minted}
\usepackage{multirow}

\title{anomaly\_vs\_outlier\_vs\_novelty}
\date{}

\begin{document}
    \maketitle
    \begin{itemize}

        \item \textbf{Outlier:}

        An outlier is an observation that deviates markedly from other members
        of the sample or appears inconsistent with the remainder of the data
        \cite{hodge-austin-2004-a-survey-of-outlier-detection-methodologies}.

        Outlierness is primarily characterized through statistical properties
        such as distance, density, probability, or deviation from the dominant
        structure of the data.

        An outlier does not necessarily represent an error or an invalid
        observation; it may result from measurement error, system faults, noise,
        or simply a rare but valid natural variation
        \cite{hodge-austin-2004-a-survey-of-outlier-detection-methodologies}.

        \item \textbf{Anomaly:}

        An anomaly is an observation, behavior, or pattern that does not conform
        to a well-defined notion of expected or normal behavior
        \cite{chandola-banerjee-kumar-2009-anomaly-detection-a-survey}.

        Depending on the application, an anomaly may correspond to a fault,
        intrusion, fraud, malfunction, or another event that is relevant to the
        system or analyst.

        Anomalies are not restricted to isolated extreme points; they may be
        point anomalies, contextual anomalies, or collective patterns whose
        abnormality emerges from their context or their relationships with other
        observations
        \cite{chandola-banerjee-kumar-2009-anomaly-detection-a-survey}.

        \item \textbf{Novelty:}

        Novelty detection is the task of identifying test observations that
        differ from the data available during training
        \cite{pimentel-clifton-clifton-tarassenko-2014-a-review-of-novelty-detection}.

        It is commonly formulated as a one-class classification problem in which a model of normality is constructed from training data that predominantly, or ideally exclusively, represent normal behavior.

        A novel observation therefore represents a previously unseen pattern relative to the learned model of normality. Unlike an anomaly, a detected
        novelty may subsequently be incorporated into the model of normal behavior if it is determined to represent a valid new regime \cite{pimentel-clifton-clifton-tarassenko-2014-a-review-of-novelty-detection,
        chandola-banerjee-kumar-2009-anomaly-detection-a-survey}.

    \end{itemize}

    \bibliographystyle{plain}
    \bibliography{/home/donkarlo/Dropbox/repo/nd_shared_research/bibliography/bibliography.bib}
\end{document}